% !PRINTED 38
\weberpage{38}
%<<CONT

thus, denoting by $p,p'$ whole rational functions,
\begin{equation}
 A'''=pA+p'A'.
\tag{21}
\end{equation}
Substituting the expressions (20), (21) in the third equation (3), one obtains for $A''''$ again an expression of the form (21), and so one may continue, obtaining finally
\begin{equation}
 A^{(\nu)}=PA+P'A',
\tag{22}
\end{equation}
where $P,P'$ are whole rational functions whose coefficients are composed from the coefficients of $A$ and $A'$ by rational operations.

In the formula (22), $A^{(\nu)}$ is a constant. This theorem is especially important in the case where $A^{(\nu)}$ is non-zero, so that $A,A'$ are relatively prime. If in this case we put
\[
 P=A^{(\nu)}F(x),
 \qquad
 P'=A^{(\nu)}\Phi(x),
\]
then, after suppressing the factor $A^{(\nu)}$, the theorem can be expressed as follows:

\gesperrt{I. If $f(x)$ and $\psi(x)$ are two whole functions without a common divisor, then two other whole functions $F(x)$ and $\Phi(x)$ can be determined which satisfy the equation}
\begin{equation}
 F(x)f(x)+\Phi(x)\psi(x)=1
\tag{23}
\end{equation}
\gesperrt{identically.}

The theorem can still be generalized. Multiplying (23) by an arbitrary whole function $\chi(x)$ gives
\begin{equation}
 F(x)\chi(x)f(x)+\Phi(x)\chi(x)\psi(x)=\chi(x),
\tag{24}
\end{equation}
and by \S~3 we can put
\begin{equation}
 \Phi(x)\chi(x)=Q(x)f(x)+\varphi(x)
\tag{25}
\end{equation}
in such a way that $Q(x),\varphi(x)$ are whole rational functions of $x$ and the degree of $\varphi(x)$ is lower than that of $f(x)$. Substituting this in (24), and putting in place of $F(x)\chi(x)+Q(x)\psi(x)$ again $F(x)$, we obtain
\[
 F(x)f(x)+\varphi(x)\psi(x)=\chi(x).
\]
We can therefore generalize the preceding theorem as follows:

\gesperrt{II. If $f(x),\psi(x),\chi(x)$ are given whole rational functions, and $f(x)$ and $\psi(x)$ are without a common divisor, then the}
%CONT>>
